% =====================================================================
%  BAB 5 — IMPLEMENTASI
% =====================================================================

\chapter{Implementasi}

\section{Lingkungan Implementasi}
\label{sec:5.1}

Implementasi dilakukan pada lingkungan yang telah dideskripsikan pada
Tabel~\ref{tab:3.lingkungan} di Bab~3. [Tambahkan detail khusus implementasi jika ada
perbedaan atau tambahan dari yang sudah disebutkan.]

\section{Implementasi [Komponen 1]}
\label{sec:5.2}

[Jelaskan bagaimana Komponen 1 diimplementasikan. Hubungkan implementasi ke rancangan
di Bab~4. Sorot keputusan teknis yang diambil selama implementasi, terutama yang berbeda
dari rancangan awal beserta alasannya.]

\subsection{[Modul / Fungsi Kunci 1]}
\label{sec:5.2.1}

[Penjelasan modul. Sertakan snippet kode untuk bagian yang non-trivial atau inovatif.
Jangan tampilkan kode yang sudah jelas; fokus pada bagian yang bernilai penjelasan.]

\begin{lstlisting}[language=Python, caption={[Deskripsi Kode Sumber 1]}, label={lst:5.kode1}]
# Contoh implementasi [fungsi kunci]
def fungsi_kunci(parameter_a, parameter_b):
    """[Docstring singkat.]"""
    # [Langkah 1: jelaskan logika non-trivial di sini]
    hasil_antara = proses_awal(parameter_a)

    # [Langkah 2]
    for item in hasil_antara:
        item = transformasi(item, parameter_b)

    return hasil_antara
\end{lstlisting}

\subsection{[Modul / Fungsi Kunci 2]}
\label{sec:5.2.2}

[Penjelasan modul berikutnya.]

\section{Implementasi [Komponen 2]}
\label{sec:5.3}

[Implementasi komponen kedua.]

\section{Implementasi [Komponen 3]}
\label{sec:5.4}

[Implementasi komponen ketiga.]

\section{Integrasi Antar-Komponen}
\label{sec:5.5}

[Jelaskan bagaimana komponen-komponen dihubungkan satu sama lain. Pipeline data, antrian
pesan, API, konfigurasi, dsb. Tampilkan diagram aliran data atau topologi deployment
jika membantu.]

\begin{figure}[H]
  \centering
  \begin{tikzpicture}[node distance=1.2cm and 2.5cm]
    % ── Ganti dengan diagram integrasi Anda ──────────────────────────
    \node[procplain, fill=cProsesBg, draw=cProses] (a) {[Komponen A]};
    \node[broker,    fill=cDataBg,   draw=cData,
          right=of a] (b) {[Broker/\\Queue]};
    \node[procplain, fill=cOkBg,    draw=cOk,
          right=of b] (c) {[Komponen C]};

    \draw[arr] (a) -- node[above,font=\scriptsize]{[data]} (b);
    \draw[arr] (b) -- node[above,font=\scriptsize]{[event]} (c);
  \end{tikzpicture}
  \caption{Topologi integrasi antar-komponen}
  \label{fig:5.integrasi}
\end{figure}

\section{Kendala Implementasi}
\label{sec:5.6}

[Uraikan kendala teknis yang ditemui selama implementasi dan bagaimana cara mengatasinya.
Ini adalah bagian yang sering dilewatkan tapi bernilai tinggi bagi pembaca.]

\begin{table}[H]
  \centering
  \caption{Kendala implementasi dan solusinya}
  \label{tab:5.kendala}
  \begin{tabular}{p{5cm}p{4cm}p{4cm}}
    \toprule
    \textbf{Kendala} & \textbf{Dampak} & \textbf{Solusi} \\
    \midrule
    [Deskripsi kendala 1] & [Dampak pada sistem] & [Solusi yang diterapkan] \\
    [Deskripsi kendala 2] & [Dampak pada sistem] & [Solusi yang diterapkan] \\
    \bottomrule
  \end{tabular}
\end{table}
